% ====================
% 6. Analysis of results, failures, limitations, costs, and risks
% ====================
\section{Analysis of Results, Failures, Limitations, Costs, and Risks}
\label{sec:discussion}

\subsection{Write-Time versus Read-Time Tagging}

Read-time tagging matches or improves adherence
relative to both memory-system baselines, whereas
write-time tagging reduces it. One possible
explanation is that write-time extraction isolates
instructions from conversational context, making
task continuity difficult to infer. Conversely,
read-time tagging retains the underlying memory
system's representation, although its tags may
be less accurate. Write-time Supermemory illustrates the
trade-off: it never over-applies a scoped memory (20/20
out-of-scope probes) but often fails to apply one when it
should.


\subsection{Why Retrieval and Tagging Are Insufficient}

The retrieval--application gap shows that access
to a relevant instruction does not ensure correct
conflict resolution. In pair \texttt{po-01}, for
example, an agent may select the standing Polars
preference despite the applicable project-specific
pandas instruction. Other failures involve
carrying a single-task exception into a new task
or combining conflicting instructions.

Effective memory thus requires both relevant information and
reliable contextual decision-making; representative examples
are in Appendix~\ref{app:failure-cases}.

\subsection{Threats to Validity}

The benchmark contains only 20 synthetic pairs with short
histories and explicit scope cues; 19 were initially
AI-assisted, and numeric conditions avoid threshold
boundaries. Findings may not generalize to longer, ambiguous,
or naturally occurring interactions. The study uses one agent
model and one run per condition; fixed sampling settings do
not guarantee reproducibility, and Supermemory's hosted
processing is not fully observable.

The scope treatments change three components at once: scope
tags, conflict-resolution rules, and memory ordering (oldest
first, versus relevance order without timestamps for the
baselines). Differences therefore cannot be attributed to
tags alone; the non-significant gain for read-time Supermemory
(37/40 versus 34/40; McNemar $p=0.25$) may arise from the
rules or ordering. Rules-only, tags-only, and ordering-only
ablations would isolate these contributions.

Retrieval matching and failure-stage diagnoses were not
independently validated in full, and the two probes within
each pair are dependent, limiting episode-level significance
tests.

\textbf{Cost.} OpenRouter spending for the full experiment was
about \$0.15 (about \$0.06 for judging); Supermemory's separately
billed processing is excluded.
